\subsubsection{机型分工与执行负载}
表\ref{tab:rev_q2_types}按正式日程重新统计运输质量、能耗和工作量。A型11架次承担231 kg，B型6架次承担140 kg，C型6架次承担387 kg。A型以较轻任务提供并行交付能力，C型以较少架次承担较大质量。机型分工由具体任务成本与资源数量共同决定，而非将需求平均拆给八架飞机。
\begin{table}[htbp]\centering\small
\caption{正式运输日程的分机型任务与负载}\label{tab:rev_q2_types}
\begin{tabular}{lrrrrrr}\toprule
机型 & 架次 & 实体数 & 质量/kg & 能耗/kWh & 累计作业/s & 平均负载/s\\\midrule
A & 11 & 4 & 231 & 26.367859 & 21118.826 & 5279.706\\
B & 6 & 2 & 140 & 15.123025 & 11133.183 & 5566.591\\
C & 6 & 2 & 387 & 24.723576 & 11360.487 & 5680.244\\
\bottomrule\end{tabular}
\end{table}


把各类累计作业时间除以实体数量，可得到忽略不可拆任务形态的平均负载下界。C型的该值为5680.243675 s，距正式最后返航仅12.987814 s；B型的平均界为5566.591421 s，其不可拆分配产生更强的离散瓶颈，见下一节。两项统计共同说明，继续改善需要改变接近满负载机型的任务结构，单纯把某一非关键任务提前并不一定改变系统工期。

在共同窗口$[0,T_{\max}]$内，定义飞机$u$的任务占用率及全机群占用率为
\begin{equation}
 U_u=\frac{\sum_{r:a_{ru}=1}p_r}{T_{\max}},\qquad
 U_{\mathrm{fleet}}=\frac{\sum_rp_r}{|\mathcal U|T_{\max}}.
 \label{eq:q2_occupation}
\end{equation}
这里的作业时间包含准备、飞行和交接，因此度量的是任务对实体资源的占用，而非纯空中飞行率。八架飞机的共同窗口占用率为95.76\%。表\ref{tab:rev_q2_util}显示，U06在该窗口上连续作业，两架C型机的占用率均超过99.7\%。这为甘特图中的紧凑安排提供了可复算的量化说明。
\begin{table}[htbp]\centering\small
\caption{共同执行时间窗口内的运输机作业占用}\label{tab:rev_q2_util}
\begin{tabular}{llrrrr}\toprule
实体 & 机型 & 架次 & 作业/min & 非占用/min & 占用率/\%\\\midrule
U01 & A & 3 & 85.808 & 9.080 & 90.43\\
U02 & A & 3 & 93.247 & 1.640 & 98.27\\
U03 & A & 3 & 88.469 & 6.418 & 93.24\\
U04 & A & 2 & 84.457 & 10.430 & 89.01\\
U05 & B & 3 & 90.666 & 4.221 & 95.55\\
U06 & B & 3 & 94.887 & 0.000 & 100.00\\
U07 & C & 3 & 94.620 & 0.268 & 99.72\\
U08 & C & 3 & 94.722 & 0.165 & 99.83\\
\bottomrule\end{tabular}
\end{table}


\subsubsection{交付进度与任务完成层次}
表\ref{tab:rev_q2_milestones}按交接完成事件统计配送进度。30 min完成28箱、268 kg，60 min完成53箱、498 kg，90 min前80箱全部交付。某一早期窗口已交付数量较少，并不表示机群没有工作：首批任务仍可能处于准备、转场或交接阶段，只有完成交接后才能进入交付累计量。
\begin{table}[htbp]\centering\small
\caption{交接完成口径下的累计配送进度}\label{tab:rev_q2_milestones}
\begin{tabular}{lrrrr}\toprule
时刻/min & 已交付箱数 & 质量/kg & 体积/m$^3$ & 箱数进度/\%\\\midrule
15 & 1 & 3 & 0.012 & 1.25\\
30 & 28 & 268 & 0.689 & 35.00\\
45 & 35 & 330 & 0.857 & 43.75\\
60 & 53 & 498 & 1.303 & 66.25\\
75 & 56 & 521 & 1.377 & 70.00\\
90 & 80 & 758 & 2.011 & 100.00\\
\bottomrule\end{tabular}
\end{table}


箱数进度和质量进度反映不同内容。医疗、食品、水和卫生用品的单箱质量不同，故某段时间交付箱数多不必然代表交付质量大。表\ref{tab:rev_delivery_kind}进一步给出各物资类别的交付范围及对应硬时限余量。实际遵守的是每只箱子的原始期限，而不是额外强加“所有医疗箱一定先于所有其他物资”的统一规则。
\begin{table}[htbp]\centering\small
\caption{各物资类别的实际交付时段与硬时限余量}\label{tab:rev_delivery_kind}
\begin{tabular}{lrrrrr}\toprule
物资类型 & 箱数 & 最早交付/min & 最后交付/min & 硬时限箱数 & 最小硬余量/min\\\midrule
医疗物资 & 16 & 14.887 & 86.220 & 16 & 3.368\\
应急食品 & 19 & 15.143 & 88.053 & 0 & --\\
饮用水 & 36 & 15.143 & 88.053 & 15 & 12.889\\
生活卫生用品 & 9 & 21.904 & 86.220 & 0 & --\\
\bottomrule\end{tabular}
\end{table}


从最后一箱交付到全部飞机返航仍有410.063102 s。这段时间属于装备回收，继续占用飞机并产生返航能耗。交付进度和最后返航分别说明物资到位与系统回收两个完成层次，使应急效果和资源周转可以同时核查。

\subsubsection{电池共享的实际接续关系}
23项任务使用14组电池，形成9条再次使用同一电池的接续关系，见表\ref{tab:rev_q2_tight_chains}。其中3条在计算精度内几乎没有充满后的等待。电池池通过不同飞机交替使用已恢复能源支撑多架次，而非一次性分配后永久绑定实体机。
\begin{table}[htbp]\centering\small
\caption{共享电池再次投入任务的接续关系}\label{tab:rev_q2_tight_chains}
\begin{tabular}{lllrrr}\toprule
电池 & 前任务 & 后任务 & 恢复/s & 后项开始/s & 接续余量/s\\\midrule
A-BAT01 & Q2-R-01 & Q2-R-15 & 2121.355 & 2178.800 & 57.445\\
A-BAT02 & Q2-R-02 & Q2-R-19 & 3600.679 & 3600.679 & 0.000\\
A-BAT03 & Q2-R-03 & Q2-R-16 & 2697.242 & 2697.242 & 0.000\\
A-BAT04 & Q2-R-04 & Q2-R-22 & 3895.818 & 3906.224 & 10.406\\
A-BAT06 & Q2-R-12 & Q2-R-23 & 4389.970 & 4389.970 & 0.000\\
B-BAT01 & Q2-R-05 & Q2-R-18 & 3012.098 & 3478.120 & 466.022\\
B-BAT02 & Q2-R-06 & Q2-R-20 & 3463.925 & 3791.400 & 327.475\\
C-BAT01 & Q2-R-07 & Q2-R-21 & 3181.215 & 3799.981 & 618.767\\
C-BAT02 & Q2-R-08 & Q2-R-17 & 2956.347 & 2989.723 & 33.376\\
\bottomrule\end{tabular}
\end{table}


对同电池的相邻任务$r,r'$，定义接续余量
\begin{equation}
 \sigma_{rr'}^B=s_{r'}-(s_r+p_r+c_r).
 \label{eq:q2_battery_slack}
\end{equation}
当恢复时间增加超过该余量，原开始时刻必须调整。接续余量较小的任务是充电压力检验的重点，但局部推迟是否增加最后返航，还取决于另一条资源链何时完成。因此后续实验沿真实依赖传播，不把所有局部延长量直接累加到工期。
